\documentclass[a4paper]{article}
\usepackage[pdftex]{graphicx}
\usepackage[utf8]{inputenc}
\usepackage{enumerate}
\usepackage{siunitx}
\sisetup{locale=DE}
\usepackage{href-ul}
\hypersetup{
	colorlinks=true,
	linkcolor=blue,
	urlcolor=blue}
\usepackage{icomma}
\usepackage{amssymb}
\usepackage{geometry}
\geometry{a4paper, top=15mm, left=15mm, right=15mm, bottom=15mm,
	headsep=10mm, footskip=12mm}

\begin{document}
	{\bf force decomposition}
	\begin{enumerate}[1.]
		
		
		{\bf \item On a ramp with the length of $l=\SI{3.0}{\meter}$, a barrel with the weight $F_G =\SI{600}{\newton}$ is brought to a $h =\SI{1,20}{\meter}$ high platform rolled.} Use a force plan to determine the required tensile force $ F_Z $.
		
		\begin{minipage}{0.6\textwidth}
			{\bf \item A truck of mass $m=\SI{28}{\tonne}$ drives uphill on a pass road with a gradient of 10\%.} After driving 20 km, it has reached the top of the pass. \\
		\end{minipage}
		\hfill
		\begin{minipage}{0.4\textwidth}
			\includegraphics[width=6 cm]{schieb115.pdf}
		\end{minipage}
		
		\begin{enumerate}[a)]
			\item Determine approximately the difference in height that the truck covers on this journey.
			\item Calculate the force with which the truck is pulled upwards by the engine during this uphill journey (if there were no friction).
		\end{enumerate}
		
		\begin{minipage}{0.5\textwidth}
			{\bf \item A sign should be attached to a house wall, as shown in the scale drawing below:}
			In order to prevent the sign from lowering due to its own weight ($F_G=\SI{52}{\newton}$), a rope should be stretched between points A and B. The horizontal pole that supports the sign is $l=\SI{2,0}{\meter}$ long ; their mass is neglected. It is attached to the wall with a hinge. The center of gravity S of the sign is at a distance of $a=\SI{1.00}{\meter}$ from the wall. The upper bracket of the rope is $h=\SI{1.50}{\meter}$ above the wall-side end of the rod.
		\end{minipage}
		\hfill
		\begin{minipage}{0.45\textwidth}
			\includegraphics[width=5.4 cm]{schild114.pdf}
		\end{minipage}
		
		Add a force plan to the drawing and determine what force the rope must absorb.
		$[\SI{1}{\centi\meter} \ \textnormal{arrow length} \ \hat{=}\ \SI{10}{\newton}]$
		
		
		\begin{minipage}{0.65\textwidth}
			{\bf \item Inclined Plane; force decomposition}\\
			Viola stands at the top of the mountain with all her ski equipment (total mass $m=\SI{51}{\kilogram}$) and bravely throws herself down a steep slope. The inclination angle $\alpha$ is $37^{\circ}$.
		\end{minipage}
		\hfill
		\begin{minipage}{0.65\textwidth}
			\includegraphics[width=4 cm]{violaski129.pdf}
		\end{minipage}
		
		\begin{enumerate}[a)]
			\item Create a force plan with the weight $\vec{F_G}$, with
			the force $\vec{F_H}$, which accelerates it towards the valley and with the force $\vec{F_N}$, which is exerted on the snow cover.\\ $[\vec{g}=\SI{9, 81}{\meter\over {\second^2}}]$
			\item How large is the acceleration $\vec{a}_0$ that Viola experiences downhill if friction forces are not taken into account?
			\item Now let the coefficient of sliding friction be $ \mu = 0.15 $. What is the acceleration now?
			$\vec{a}_R$?
		\end{enumerate}
		
	\end{enumerate}
	
	\fbox{
		\begin{minipage}{0.5\textwidth}
			For the solution, please \href{https://www.okuyakl.de/phys/p9kraL404/le404.pdf}{click here} or scan the QR code.\\
			Additional worksheets are available at
			
			\href{http://www.okuyakl.com}{www.okuyakl.com}
		\end{minipage}
		\hfill
		\begin{minipage}{0.4\textwidth}
			\includegraphics[width=1.5 cm]{../../viecher/zwe03}
			\includegraphics[width=3 cm]{qre404}
			\includegraphics[width=2 cm]{../../viecher/afanticon1}
			
	\end{minipage}}
	
\end{document}